\section{Test plan}\label{sec:10}

\subsection{How the system is verified}\label{sec:10.1}

Three levels, and the traceability rule that every test names the requirement ID it verifies.
Unit tests live beside the code and cover the \textbf{control} and \textbf{entity} layers, which is
where the business logic is: the value objects and the monitor state machine in \texttt{shared/},
the mass balance, the flow split and fault injection in \texttt{physics-simulator/}, the hysteresis
loop in \texttt{dcv-controller/}, the CUSUM and the classifier in \texttt{fault-monitor/}, the
compensation solver in \texttt{supervisor/}, the view projection in \texttt{dashboard/}. Because
the clock is a port, the \textbf{offline harness} in \texttt{evaluation/} drives the
same control and diagnosis classes the containers run against a simulated clock, so a seven-day
replay and twenty injection runs cost seconds rather than days. End-to-end and fault-injection runs
go through the real compose stack.

\begin{longtable}[]{@{}
  >{\raggedright\arraybackslash}p{(\columnwidth - 6\tabcolsep) * \real{0.07}}
  >{\raggedright\arraybackslash}p{(\columnwidth - 6\tabcolsep) * \real{0.33}}
  >{\raggedright\arraybackslash}p{(\columnwidth - 6\tabcolsep) * \real{0.14}}
  >{\raggedright\arraybackslash}p{(\columnwidth - 6\tabcolsep) * \real{0.46}}@{}}
\toprule
\textbf{ID} & \textbf{Test or experiment} & \textbf{Requirement} & \textbf{Has it run, and what came
out} \\
\midrule
\endhead
T-24 & Detection delay swept against the CUSUM limit $h$, ≥ 20 runs per point & R-1 &
\textbf{Run}, 200 runs over six arms: stepped p95 6 min, 20/20; closed-loop p95 257 min, 15/20 \\
T-25 & Model-based monitor against a plain residual threshold, identical seeds, 7-day fault-free
replay & R-1, R-3 & \textbf{Run}: at equal alarm volume the monitor detects 0.75 against the
threshold's 0.15 at a median delay of 74 min against 312; R-3 \textbf{fails} at 0.55 confirmed cases
per fault-free week, worst week 6 \\
T-15 & Confusion matrix over the four injected faults on a fixed seed, plus a no-fault arm &
R-1, R-3 & \textbf{Run}, 218 runs: \texttt{STUCK\_DAMPER} named correctly,
\texttt{BLOCKED\_FILTER} → \texttt{UNCLASSIFIED} 11/11, \texttt{SENSOR\_DRIFT} \textbf{never
emitted}; \textbf{one of four mechanisms classified} \\
T-26 & Comfort during a latched fault, with and without compensation & R-2 & \textbf{Run}: 5 of 5
seeds on the drawn weekday, \textbf{0 of 5 at design occupancy}, no-fault reference 5 of 5
throughout \\
T-27 & Kill each container in turn & R-4 & \textbf{Run} in-JVM, 23 of 23 assertions:
\texttt{BLIND} in 4 cycles, 0 fault events, 0 commands on stale data; container kills \textbf{now
also run} against \texttt{docker compose}, five of six rooms \texttt{BLIND} in 12 s \\
T-20 & Rebuild a random diagnosis from the journal & R-5 & \textbf{Demonstrated live} in the
recorded demo run: HTTP 200, 77\,248 bytes, \textbf{445 events}; the id was the one that run produced
rather than a drawn one. \textbf{And now tested} by eight integration tests on the real SQL \\
T-28 & Scale zones per monitor to the breaking point & R-6 & \textbf{Run}: 25\,600 zones at a p95
of about 6 s measured, $\approx$ 86\,000 \textbf{extrapolated},
\texttt{AhuFlowModel.believedFlowTo} named \\
T-34b & Room isolation: one injected fault per run, rotating the injected room over ≥ 3 rooms &
R-7 & \textbf{Run}, both occupancy regimes, 540 observed rooms. Drawn weekday \textbf{0 of 270
failed}; design occupancy \textbf{231 of 270 failed} \\
T-29 & Optional learned one-class contender against the CUSUM on identical seeds & R-1, R-3 &
\textbf{Run}: \textbf{dropped}. To match the monitor's 0.75 detection it costs about 20× the false
alarms (11.25 against 0.55 per week) and 7× the delay (518 against 74 min) \\
T-35 & Classify after a probe has been judged, or at confirmation? 176 paired runs on identical
seeds & R-1 & \textbf{Run}: \textbf{not adopted}. Detection moves in 0 of 88 pairs and the settled
classification is identical in all 88, because the judged step reads \texttt{MISSING} for every
condition \\
T-36 & One observation published on the real broker, consumed by the running
\texttt{fault-monitor} and asserted in the real journal & R-5 & \textbf{Run}, 2 tests against the
deployed stack: the row carries the published wattage, the publisher's instant and the monitor's
own provenance; a reading for a room outside the zone set journals nothing. \textbf{Skips} with the
reason when the stack is down \\
T-37 & A latched diagnosis survives a \texttt{SIGKILL} and restart of its container & --- &
\textbf{Run live on the demo profile}, transcript in
\texttt{evaluation/\allowbreak results/\allowbreak T-37-\allowbreak restart-\allowbreak recovery.txt}. A115 confirmed \texttt{STUCK\_DAMPER};
\texttt{docker kill} and restart; the case id and \texttt{confirmedAt} are \emph{identical}
across the kill while \texttt{updatedAt} keeps advancing. A cold start would have minted both
afresh. The damper pin and the believed fan speed come back with it \\
\bottomrule
\end{longtable}

\textbf{Appendix \ref{sec:K}} carries the same eleven rows with the expected result each was written against
and the evidence in full; §\ref{sec:11.2} and Appendix \ref{sec:F} report the experiments themselves.

Four named unit tests are worth quoting, because they are what the strongest claims in this report
actually rest on. Test names describe the behaviour and carry the requirement ID, so the
traceability runs all the way into the suite:

\begin{footnotesize}
\begin{verbatim}
shouldHoldEveryHealthyFlowExactlyWhileTheFanHasHeadroom_R7    zero healthy deviation,
                                                              gain exactly 1/t
shouldRestoreTheFrozenRoomToItsDesignFlowAtTheNormalFanSpeed_R2   59.7 -> 165.9 l/s
shouldSpendExactlyTheStatedToleranceWhenTheFanIsAlreadyAtNominal_R7   the saturated fan
shouldRefuseToCompensateADamperWithNoAirPath_R2Bound          no command at all
\end{verbatim}
\end{footnotesize}

\subsection{The state of the suite, stated plainly}\label{sec:10.2}

A full \texttt{mvn -o clean test} run on \factSuiteRunDate{} reports \textbf{BUILD SUCCESS over all
\factReactorModules{} reactor modules: \factSuiteTests{} tests, 0 failures, 0 errors}.
\textbf{The skip count is not a fixed property of the suite, and no skip is a gap}:
\factGuardedTests{} tests assume something outside the JVM is reachable, the
\factJournalTests{} journal tests that execute R-5's real SQL against a PostgreSQL (T-20), the
\factPipelineTests{} pipeline tests that need the broker and a running \texttt{fault-monitor} as
well (T-36), and the \factReconnectTests{} reconnect test that needs the broker, so they skip
when nothing is up and run when the stack is. The run above was taken with the stack up, which is
why it reports \textbf{0 skipped} rather than \factGuardedTests{}. Starting anything is
deliberately not the default, so the suite stays runnable in seconds with nothing installed.

What was the single largest gap in this project is closed: \texttt{evaluation/} holds its own sources
(a \texttt{Main} and \textbf{\factExperiments{}} experiments), and a full run on OpenJDK 25.0.3 completes in
\textbf{350 s} of wall clock, writing the nine result tables §\ref{sec:11.2} reports from. That run is also the
evidence for §\ref{sec:13.9}'s determinism claim, and it is worth stating as a result: re-running the harness
against the committed result files leaves \textbf{eight of the nine byte-identical}. The ninth is
T-28, and it does not reproduce because it is the one experiment that measures \emph{wall clock}
rather than simulated behaviour, so the exception is the experiment where bit-for-bit reproduction
would mean the measurement was not being taken. \textbf{Appendix \ref{sec:K}} gives the per-module test counts,
the per-experiment timings and both results in full.
